\section{introduction}
Avant d’entamer la phase de conception et de développement de notre application web, il est essentiel de comprendre le contexte dans lequel ce projet s’inscrit. L’étude préalable permet d’identifier les besoins réels, les difficultés rencontrées dans le processus actuel, ainsi que les objectifs à atteindre.
Ce chapitre pose les bases du projet en abordant la problématique, la solution proposée, ainsi que les besoins fonctionnels et non fonctionnels nécessaires au bon fonctionnement de l’application.
\section{problemeatique}
Dans de nombreuses entreprises industrielles, la gestion de la maintenance et la supervision de la production restent encore réalisées de manière manuelle, dispersée ou à l’aide d’outils non centralisés.
Cela entraîne plusieurs difficultés :

Un manque de visibilité sur l’état réel des équipements.

Une difficulté à planifier les interventions de maintenance préventive.

Une perte de temps lors de la recherche d'informations ou de l’historique des pannes.

Une absence d’indicateurs fiables pour évaluer la performance des machines et de la production.

Des erreurs humaines dues à la saisie manuelle.

Face à ces limitations, il devient nécessaire de mettre en place une solution numérique qui facilite la gestion, améliore la coordination entre les équipes et optimise le suivi global des activités.
\section{solition proposee}
\section{besoin fonctionnel et non fonctionnel}
\section{conclusion}